\documentclass[14pt,a4paper]{extarticle}

\usepackage[utf8]{inputenc}
\usepackage[T2A]{fontenc}
\usepackage[russian]{babel}
\usepackage{pscyr}

\usepackage{geometry}
\geometry{left=30mm,right=15mm,top=20mm,bottom=20mm}

\usepackage{setspace}
\singlespacing
\renewcommand{\baselinestretch}{1.0}
\setlength{\parskip}{0pt}

\usepackage{indentfirst}
\setlength{\parindent}{1.25cm}

\usepackage{amsmath,amssymb}
\usepackage{graphicx}
\usepackage{float}
\usepackage{caption}
\usepackage{enumitem}
\usepackage{booktabs}
\usepackage{tabularx}
\usepackage{array}

\usepackage{titlesec}

\titleformat{\section}
  {\normalfont\fontsize{14}{18}\bfseries\filright}
  {\thesection}{1em}{}
\titlespacing*{\section}{1.25cm}{14pt}{14pt}

\titleformat{\subsection}
  {\normalfont\fontsize{14}{18}\bfseries\filright}
  {\thesubsection}{1em}{}
\titlespacing*{\subsection}{1.25cm}{14pt}{14pt}

\usepackage{tocloft}
\renewcommand{\cfttoctitlefont}{\hfill\fontsize{14}{18}\selectfont\bfseries}
\renewcommand{\cftaftertoctitle}{\hfill}
\renewcommand{\cftsecfont}{\normalfont}
\renewcommand{\cftsecpagefont}{\normalfont}
\renewcommand{\cftsubsecfont}{\normalfont}
\renewcommand{\cftsubsecpagefont}{\normalfont}
\renewcommand{\cftdotsep}{1}
\renewcommand{\cftsecleader}{\cftdotfill{\cftdotsep}}

\usepackage{listings}
\lstset{
  basicstyle=\ttfamily\small,
  frame=single,
  breaklines=true,
  showstringspaces=false,
  inputencoding=utf8,
  extendedchars=true,
  literate={а}{{\cyra}}1 {б}{{\cyrb}}1 {в}{{\cyrv}}1 {г}{{\cyrg}}1
           {д}{{\cyrd}}1 {е}{{\cyre}}1 {ж}{{\cyrzh}}1 {з}{{\cyrz}}1
           {и}{{\cyri}}1 {й}{{\cyrishrt}}1 {к}{{\cyrk}}1 {л}{{\cyrl}}1
           {м}{{\cyrm}}1 {н}{{\cyrn}}1 {о}{{\cyro}}1 {п}{{\cyrp}}1
           {р}{{\cyrr}}1 {с}{{\cyrs}}1 {т}{{\cyrt}}1 {у}{{\cyru}}1
           {ф}{{\cyrf}}1 {х}{{\cyrh}}1 {ц}{{\cyrc}}1 {ч}{{\cyrch}}1
           {ш}{{\cyrsh}}1 {щ}{{\cyrshch}}1 {ъ}{{\cyrhrdsn}}1 {ы}{{\cyrery}}1
           {ь}{{\cyrsftsn}}1 {э}{{\cyrerev}}1 {ю}{{\cyryu}}1 {я}{{\cyrya}}1
           {А}{{\CYRA}}1 {Б}{{\CYRB}}1 {В}{{\CYRV}}1 {Г}{{\CYRG}}1
           {Д}{{\CYRD}}1 {Е}{{\CYRE}}1 {Ж}{{\CYRZH}}1 {З}{{\CYRZ}}1
           {И}{{\CYRI}}1 {Й}{{\CYRISHRT}}1 {К}{{\CYRK}}1 {Л}{{\CYRL}}1
           {М}{{\CYRM}}1 {Н}{{\CYRN}}1 {О}{{\CYRO}}1 {П}{{\CYRP}}1
           {Р}{{\CYRR}}1 {С}{{\CYRS}}1 {Т}{{\CYRT}}1 {У}{{\CYRU}}1
           {Ф}{{\CYRF}}1 {Х}{{\CYRH}}1 {Ц}{{\CYRC}}1 {Ч}{{\CYRCH}}1
           {Ш}{{\CYRSH}}1 {Щ}{{\CYRSHCH}}1 {Ъ}{{\CYRHRDSN}}1 {Ы}{{\CYRERY}}1
           {Ь}{{\CYRSFTSN}}1 {Э}{{\CYREREV}}1 {Ю}{{\CYRYU}}1 {Я}{{\CYRYA}}1
           {ё}{{\cyryo}}1 {Ё}{{\CYRYO}}1,
  numbers=left,
  numberstyle=\tiny,
  tabsize=2,
  language=R
}

\usepackage{hyperref}
\hypersetup{
    colorlinks=true,
    linkcolor=black,
    urlcolor=blue,
    citecolor=black,
}

\usepackage{fancyhdr}
\pagestyle{fancy}
\fancyhf{}
\fancyfoot[R]{\thepage}
\renewcommand{\headrulewidth}{0pt}

\newcolumntype{Y}{>{\raggedright\arraybackslash}X}
\sloppy

\begin{document}

% =========================================================
% ТИТУЛЬНЫЙ ЛИСТ
% =========================================================
\thispagestyle{empty}
\begin{center}
МИНИСТЕРСТВО ОБРАЗОВАНИЯ РЕСПУБЛИКИ БЕЛАРУСЬ

\vspace{0.5em}

Учреждение образования

\textbf{<<Белорусский государственный университет\\информатики и радиоэлектроники>>}

\vspace{2em}

Факультет компьютерных систем и сетей

\vspace{0.5em}

Кафедра программного обеспечения информационных технологий
\end{center}

\vspace{4em}

\begin{center}
\textbf{ОТЧЁТ-ИНСТРУКЦИЯ}

\vspace{0.5em}

по лабораторным работам №1 и №2

\vspace{0.5em}

по дисциплине <<Стили и методы программирования>>

\vspace{1em}

\textbf{Тема: <<Инструкция по установке, настройке и эксплуатации RStudio. Подготовка и выполнение скриптов на языке R>>}
\end{center}

\vspace{5em}

\begin{flushright}
\begin{minipage}{0.5\textwidth}
\textbf{Выполнил:}\\
Лебедевич Артем Владимирович\\
магистрант факультета\\
компьютерных систем и сетей\\
группы 556301

\vspace{1em}

\textbf{Преподаватель:}\\
Парамонов Антон Иванович\\
кандидат технических наук, доцент\\
заведующий кафедрой информационных\\
систем и технологий ИТТ БГУИР
\end{minipage}
\end{flushright}

\vfill

\begin{center}
Минск, 2026
\end{center}

\newpage

% =========================================================
% СОДЕРЖАНИЕ
% =========================================================
\tableofcontents
\newpage

% =========================================================
% ВВЕДЕНИЕ
% =========================================================
\addcontentsline{toc}{section}{Введение}
\section*{ВВЕДЕНИЕ}

Данный документ подготовлен как практическая инструкция к лабораторным работам №1 и №2 по языку R. Его задача не просто повторить листинг программы, а объяснить, что именно делается на каждом шаге, зачем используется тот или иной метод, какие параметры ему передаются и как эти параметры влияют на результат.

Лабораторная работа №1 посвящена знакомству с R и средой RStudio: созданию проекта, запуску скриптов, работе с векторами, матрицами, фреймами данных и базовыми графиками. Лабораторная работа №2 продолжает эту тему на более прикладном примере: данные загружаются из удалённого API, сохраняются в CSV-файл, затем анализируются и визуализируются с помощью внешних библиотек.

Инструкция написана с расчётом на студента, который только открывает RStudio и хочет понимать не только <<что нажать>>, но и <<почему именно так>>. Поэтому рядом с командами приводятся пояснения по вкладкам интерфейса, библиотекам, параметрам функций и типовым ошибкам.

\newpage

% =========================================================
% РАЗДЕЛ 1
% =========================================================
\section{ИНСТРУКЦИЯ ПО УСТАНОВКЕ, НАСТРОЙКЕ И ЭКСПЛУАТАЦИИ RSTUDIO}

Для выполнения лабораторных работ используется связка из двух программных компонентов: языка R и среды разработки RStudio. Важно понимать разницу между ними. R является самим языком и интерпретатором, который выполняет команды. RStudio является удобной оболочкой вокруг R: в ней можно писать скрипты, запускать их по частям, смотреть переменные, строить графики и работать с файлами проекта.

\subsection{Установка R}

Сначала устанавливается язык R. Для этого необходимо открыть официальный сайт CRAN: \url{https://cran.r-project.org/}. На сайте выбирается версия для своей операционной системы: Windows, macOS или Linux.

Рекомендуемый порядок установки:

\begin{enumerate}[nosep]
  \item Скачать установщик R с CRAN.
  \item Запустить установщик и оставить стандартные параметры, если нет специальных требований.
  \item После установки проверить, что R запускается: открыть терминал или консоль R и выполнить команду \texttt{R.version.string}.
\end{enumerate}

Главный смысл этого шага: RStudio не заменяет R. Если установить только RStudio, но не установить сам R, среда разработки не сможет выполнять скрипты.

\subsection{Установка RStudio}

После установки R загружается RStudio Desktop с сайта Posit: \url{https://posit.co/downloads/}. Для лабораторных работ достаточно бесплатной версии RStudio Desktop.

На macOS приложение обычно перетаскивается в папку \texttt{Applications}. На Windows используется обычный мастер установки. После первого запуска RStudio автоматически находит установленный интерпретатор R. Если на компьютере установлено несколько версий R, нужную версию можно выбрать в настройках RStudio.

\subsection{Создание проекта}

Проект в RStudio нужен для того, чтобы все скрипты, данные и результаты лежали в одной рабочей папке. Это снижает количество ошибок с путями к файлам.

Порядок создания проекта:

\begin{enumerate}[nosep]
  \item Открыть меню \texttt{File -> New Project...}.
  \item Выбрать \texttt{New Directory -> New Project}.
  \item Указать имя папки проекта и место её расположения.
  \item Нажать \texttt{Create Project}.
\end{enumerate}

После этого RStudio создаёт файл с расширением \texttt{.Rproj}. Его удобно открывать в следующий раз двойным щелчком: RStudio сразу восстановит рабочую папку проекта.

\subsection{Подготовка и запуск R-скрипта}

Скрипт на языке R сохраняется в файле с расширением \texttt{.R}. Для лабораторных работ используются файлы \texttt{lab1\_script.R}, \texttt{lab2\_script1.R} и \texttt{lab2\_script2.R}.

Основные способы запуска:

\begin{enumerate}[nosep]
  \item \texttt{Ctrl+Enter} или \texttt{Cmd+Enter} на macOS: выполнить текущую строку или выделенный фрагмент. Этот способ удобен, когда надо пройти лабораторную пошагово.
  \item Кнопка \texttt{Source}: выполнить весь скрипт целиком. Этот способ удобен после проверки отдельных строк.
  \item Команда \texttt{source("lab1\_script.R")} в консоли: запустить файл как программу.
  \item Команда \texttt{Rscript lab1\_script.R} из терминала: выполнить скрипт без открытия RStudio.
\end{enumerate}

Для учебной работы лучше начинать с первого варианта. Тогда видно, какая команда что создаёт, какой график появляется и на каком шаге возникла ошибка.

\newpage

% =========================================================
% РАЗДЕЛ 2
% =========================================================
\section{ИНТЕРФЕЙС RSTUDIO: ЧТО ГДЕ НАХОДИТСЯ И ЗАЧЕМ НУЖНО}

Окно RStudio обычно разделено на четыре рабочие области. На первых занятиях кажется, что вкладок слишком много, но каждая из них решает конкретную задачу.

\begin{table}[H]
\caption{Основные области и вкладки RStudio}
\begin{tabularx}{\textwidth}{|l|Y|Y|}
\hline
\textbf{Вкладка или область} & \textbf{Что показывает} & \textbf{Что можно сделать при выполнении ЛР} \\
\hline
\texttt{Source} & Текст открытого R-скрипта. & Написать код, выделить нужный фрагмент, выполнить строку через \texttt{Ctrl+Enter}, сохранить файл. \\
\hline
\texttt{Console} & Команды R и результаты их выполнения. & Смотреть вывод \texttt{print()}, \texttt{cat()}, \texttt{summary()}, сообщения об ошибках и предупреждения. \\
\hline
\texttt{Environment} & Созданные переменные, векторы, таблицы, модели. & Проверить, появились ли \texttt{data1}, \texttt{A}, \texttt{df}, \texttt{users\_data}, \texttt{kmeans\_result}. \\
\hline
\texttt{History} & История ранее выполненных команд. & Быстро повторить команду, если она уже запускалась вручную. \\
\hline
\texttt{Files} & Файлы и папки текущего проекта. & Проверить наличие скриптов, CSV-файла и сохранённых PNG-графиков. \\
\hline
\texttt{Plots} & Графики, построенные командами R. & Посмотреть результат \texttt{barplot()}, \texttt{mosaicplot()}, \texttt{ggplot()}, экспортировать график. \\
\hline
\texttt{Packages} & Установленные и подключённые пакеты. & Проверить наличие \texttt{httr}, \texttt{jsonlite}, \texttt{ggplot2}, \texttt{dplyr}. \\
\hline
\texttt{Help} & Справка по функциям и пакетам. & Открыть документацию, например \texttt{?barplot}, \texttt{?kmeans}, \texttt{?ggsave}. \\
\hline
\texttt{Viewer} & HTML-результаты, интерактивные отчёты и web-содержимое. & В этих лабораторных используется редко, но полезен при работе с HTML-виджетами. \\
\hline
\end{tabularx}
\end{table}

При выполнении ЛР 1 особенно важны вкладки \texttt{Source}, \texttt{Console}, \texttt{Environment} и \texttt{Plots}. Например, после команды создания вектора \texttt{data1} он появляется в \texttt{Environment}; после \texttt{barplot()} график появляется в \texttt{Plots}; после \texttt{summary(df)} результат выводится в \texttt{Console}.

Удобная логика работы такая: пишем или открываем код в \texttt{Source}, запускаем его по частям, смотрим текстовый результат в \texttt{Console}, проверяем переменные в \texttt{Environment}, графики смотрим в \texttt{Plots}. Если файл не найден, первым делом проверяется вкладка \texttt{Files} и текущая рабочая папка.

\newpage

% =========================================================
% РАЗДЕЛ 3
% =========================================================
\section{ЛАБОРАТОРНАЯ РАБОТА №1: ПОШАГОВОЕ ВЫПОЛНЕНИЕ}

В ЛР 1 используется базовый R без подключения внешних библиотек. Это важно: все методы из данного раздела входят в стандартную поставку языка R, поэтому отдельные команды \texttt{install.packages()} и \texttt{library()} для них не нужны.

\subsection{Шаг 1. Создание числового вектора}

Сначала создаётся вектор оценок:

\begin{lstlisting}
data1 <- c(8, 5, 9, 7, 3, 10, 6, 4)
\end{lstlisting}

Метод \texttt{c()} используется для объединения значений в один вектор. В R вектор является одной из базовых структур данных. Здесь вектор хранит оценки студентов.

\begin{table}[H]
\caption{Метод \texttt{c()} в ЛР 1}
\begin{tabularx}{\textwidth}{|l|Y|}
\hline
\textbf{Параметр} & \textbf{Назначение и влияние} \\
\hline
Набор значений & Значения перечисляются через запятую. Порядок важен: первый элемент вектора считается первым студентом, второй элемент -- вторым и так далее. \\
\hline
\end{tabularx}
\end{table}

\subsection{Шаг 2. Назначение имён элементам вектора}

Далее элементам вектора присваиваются фамилии:

\begin{lstlisting}
names(data1) <- c("Иванов", "Петров", "Сидоров", "Козлов",
                  "Новиков", "Фёдоров", "Морозов", "Волков")
\end{lstlisting}

Метод \texttt{names()} используется для задания имён элементам объекта. После этого значения \texttt{data1} можно воспринимать не как безымянные числа, а как оценки конкретных студентов.

\begin{table}[H]
\caption{Метод \texttt{names()}}
\begin{tabularx}{\textwidth}{|l|Y|}
\hline
\textbf{Что указывается} & \textbf{Как влияет на работу} \\
\hline
\texttt{names(data1)} & Указывает, какому объекту назначаются имена. \\
\hline
Вектор строк & Количество имён должно соответствовать количеству элементов \texttt{data1}. Если имён меньше или больше, данные становятся неудобными для чтения и могут появиться \texttt{NA}. \\
\hline
\end{tabularx}
\end{table}

\subsection{Шаг 3. Построение столбиковой диаграммы}

Для визуализации оценок используется встроенный метод \texttt{barplot()}:

\begin{lstlisting}
barplot(data1,
        main = "Успеваемость студентов",
        col = rainbow(length(data1)),
        ylab = "Оценка",
        las = 2,
        cex.names = 0.8)
\end{lstlisting}

Метод \texttt{barplot()} строит столбиковую диаграмму. В данной лабораторной он нужен, чтобы быстро увидеть, у кого оценка выше, а у кого ниже.

\begin{table}[H]
\caption{Основные параметры \texttt{barplot()}}
\begin{tabularx}{\textwidth}{|l|Y|}
\hline
\textbf{Параметр} & \textbf{Назначение и влияние} \\
\hline
\texttt{height} или первый аргумент & Набор чисел, по которым строятся столбцы. В примере это \texttt{data1}. \\
\hline
\texttt{main} & Заголовок графика. Помогает понять, что изображено. \\
\hline
\texttt{col} & Цвета столбцов. В примере \texttt{rainbow(length(data1))} создаёт столько цветов, сколько элементов в векторе. \\
\hline
\texttt{ylab} & Подпись вертикальной оси. Здесь ось показывает оценку. \\
\hline
\texttt{las = 2} & Поворачивает подписи оси, чтобы фамилии лучше помещались под столбцами. \\
\hline
\texttt{cex.names = 0.8} & Уменьшает размер подписей. Чем меньше значение, тем компактнее текст. \\
\hline
\end{tabularx}
\end{table}

Дополнительно применяется метод \texttt{rainbow()}. Он относится к базовой графике R и создаёт палитру цветов. Параметр \texttt{length(data1)} означает, что цветов должно быть ровно столько, сколько столбцов.

\subsection{Шаг 4. Создание строкового вектора}

Во втором векторе хранятся имена студентов:

\begin{lstlisting}
data2 <- c("Алексей", "Борис", "Виктор", "Григорий",
           "Дмитрий", "Евгений", "Жан", "Захар")
\end{lstlisting}

Здесь снова используется метод \texttt{c()}, но теперь элементы являются строками. Для R это обычный вектор, только тип данных не числовой, а символьный.

\subsection{Шаг 5. Сортировка оценок и согласование имён}

Чтобы показать оценки от большей к меньшей, применяется сортировка:

\begin{lstlisting}
sorted_data1 <- sort(data1, decreasing = TRUE)
sorted_data2 <- data2[order(data1, decreasing = TRUE)]
\end{lstlisting}

Метод \texttt{sort()} возвращает отсортированные значения. Метод \texttt{order()} возвращает порядок индексов, по которому надо переставить другой вектор. В данной работе это важно: если отсортировать только оценки, имена студентов могут перестать соответствовать своим оценкам. Поэтому \texttt{order()} используется для перестановки \texttt{data2} тем же образом.

\begin{table}[H]
\caption{Методы сортировки}
\begin{tabularx}{\textwidth}{|l|Y|Y|}
\hline
\textbf{Метод} & \textbf{Зачем используется} & \textbf{Важные параметры} \\
\hline
\texttt{sort()} & Получить значения в отсортированном виде. & \texttt{decreasing = TRUE} сортирует по убыванию, \texttt{FALSE} -- по возрастанию. \\
\hline
\texttt{order()} & Получить индексы элементов в нужном порядке. & Используется с тем же параметром \texttt{decreasing}, чтобы связанные данные переставлялись одинаково. \\
\hline
\end{tabularx}
\end{table}

\subsection{Шаг 6. Добавление легенды к графику}

После построения отсортированной диаграммы добавляется легенда:

\begin{lstlisting}
legend("topright",
       legend = sorted_data2,
       fill = rainbow(length(sorted_data2)),
       cex = 0.7,
       title = "Студенты")
\end{lstlisting}

Метод \texttt{legend()} нужен, чтобы расшифровать цвета или обозначения на графике.

\begin{table}[H]
\caption{Параметры \texttt{legend()}}
\begin{tabularx}{\textwidth}{|l|Y|}
\hline
\textbf{Параметр} & \textbf{Назначение и влияние} \\
\hline
\texttt{"topright"} & Положение легенды. Можно использовать также \texttt{"topleft"}, \texttt{"bottomright"} и другие варианты. \\
\hline
\texttt{legend} & Текстовые подписи внутри легенды. В примере это имена студентов. \\
\hline
\texttt{fill} & Цветные прямоугольники рядом с подписями. Цвета должны соответствовать цветам столбцов. \\
\hline
\texttt{cex} & Размер текста легенды. Значение \texttt{0.7} делает подписи немного меньше. \\
\hline
\texttt{title} & Заголовок легенды. \\
\hline
\end{tabularx}
\end{table}

\subsection{Шаг 7. Создание матрицы}

Матрица создаётся методом \texttt{matrix()}:

\begin{lstlisting}
A <- matrix(c(2, 5, 8, 1,
              3, 7, 4, 6,
              9, 2, 5, 3,
              1, 8, 6, 4,
              7, 3, 2, 9),
            nrow = 4, ncol = 5)
\end{lstlisting}

Метод \texttt{matrix()} преобразует набор значений в таблицу фиксированного размера. В лабораторной создаётся матрица размером 4 на 5.

\begin{table}[H]
\caption{Параметры \texttt{matrix()}}
\begin{tabularx}{\textwidth}{|l|Y|}
\hline
\textbf{Параметр} & \textbf{Назначение и влияние} \\
\hline
Первый аргумент & Значения, из которых заполняется матрица. \\
\hline
\texttt{nrow = 4} & Количество строк. В примере это 4 группы. \\
\hline
\texttt{ncol = 5} & Количество столбцов. В примере это 5 предметов. \\
\hline
\texttt{byrow} & Если указать \texttt{byrow = TRUE}, матрица будет заполняться по строкам. По умолчанию R заполняет матрицу по столбцам. \\
\hline
\end{tabularx}
\end{table}

Названия строк и столбцов задаются методами \texttt{rownames()} и \texttt{colnames()}. Это делает матрицу понятнее при выводе и при построении графика.

\subsection{Шаг 8. Построение мозаичной диаграммы}

Для матрицы строится мозаичная диаграмма:

\begin{lstlisting}
mosaicplot(A,
           main = "Мозаичная диаграмма матрицы A",
           color = heat.colors(ncol(A)),
           las = 1)
\end{lstlisting}

Метод \texttt{mosaicplot()} показывает соотношения между категориями в виде прямоугольников. Чем больше значение в ячейке матрицы, тем заметнее соответствующая часть диаграммы.

\begin{table}[H]
\caption{Параметры \texttt{mosaicplot()}}
\begin{tabularx}{\textwidth}{|l|Y|}
\hline
\textbf{Параметр} & \textbf{Назначение и влияние} \\
\hline
Первый аргумент & Таблица или матрица, по которой строится диаграмма. \\
\hline
\texttt{main} & Заголовок диаграммы. \\
\hline
\texttt{color} & Цветовая схема. В примере \texttt{heat.colors(ncol(A))} создаёт тёплую палитру по числу столбцов. \\
\hline
\texttt{las = 1} & Настраивает ориентацию подписей осей, чтобы они легче читались. \\
\hline
\end{tabularx}
\end{table}

\subsection{Шаг 9. Создание фрейма данных и итоговая статистика}

В конце данные объединяются в таблицу:

\begin{lstlisting}
df <- data.frame(Имя = data2, Оценка = as.numeric(data1))
print(df)
summary(df)
\end{lstlisting}

Метод \texttt{data.frame()} создаёт табличную структуру, похожую на таблицу Excel: каждый столбец имеет имя и свой тип данных. Метод \texttt{as.numeric()} нужен, чтобы гарантировать числовой тип оценок. Метод \texttt{summary()} выводит краткую статистику: минимум, максимум, среднее, медиану и квартильные значения для числовых столбцов.

\begin{table}[H]
\caption{Итоговые методы ЛР 1}
\begin{tabularx}{\textwidth}{|l|Y|Y|}
\hline
\textbf{Метод} & \textbf{Назначение} & \textbf{На что обратить внимание} \\
\hline
\texttt{data.frame()} & Создаёт таблицу из нескольких векторов. & Длина всех столбцов должна совпадать. \\
\hline
\texttt{as.numeric()} & Преобразует данные к числовому типу. & Нужно для корректных расчётов и статистики. \\
\hline
\texttt{print()} & Выводит объект в консоль. & Удобно для проверки результата. \\
\hline
\texttt{summary()} & Показывает краткое описание данных. & Для чисел выводит статистику, для строк -- другую сводную информацию. \\
\hline
\end{tabularx}
\end{table}

\newpage

% =========================================================
% РАЗДЕЛ 4
% =========================================================
\section{МЕТОД УКАЗАНИЯ ПОДГОТОВКИ СКРИПТА ДЛЯ ВЫПОЛНЕНИЯ ЛР 2}

Лабораторная работа №2 состоит из двух логических частей. Первый скрипт подключается к удалённому ресурсу и сохраняет данные. Второй скрипт читает сохранённый CSV-файл, выполняет анализ и строит графики.

Такое разделение удобно: если данные уже скачаны, не нужно каждый раз обращаться к API. Можно отдельно отлаживать анализ и визуализацию.

\subsection{Какие библиотеки нужны}

В ЛР 2 используются внешние библиотеки R. Если они ещё не установлены, их надо установить один раз:

\begin{lstlisting}
install.packages(c("httr", "jsonlite", "ggplot2", "dplyr", "corrplot", "cluster"))
\end{lstlisting}

После установки в каждом новом сеансе R нужные библиотеки подключаются командой \texttt{library()}.

\begin{table}[H]
\caption{Библиотеки для ЛР 2}
\begin{tabularx}{\textwidth}{|l|Y|Y|}
\hline
\textbf{Библиотека} & \textbf{Для чего используется} & \textbf{Где применяется} \\
\hline
\texttt{httr} & Выполнение HTTP-запросов к удалённым ресурсам. & \texttt{GET()}, \texttt{status\_code()}, \texttt{content()} в первом скрипте. \\
\hline
\texttt{jsonlite} & Преобразование JSON-ответа в таблицу R. & \texttt{fromJSON()} в первом скрипте. \\
\hline
\texttt{ggplot2} & Построение графиков по грамматике графики. & \texttt{ggplot()}, \texttt{geom\_point()}, \texttt{geom\_bar()}, \texttt{ggsave()} во втором скрипте. \\
\hline
\texttt{dplyr} & Удобная обработка таблиц. & \texttt{select()}, \texttt{mutate()} во втором скрипте. \\
\hline
\texttt{corrplot} & Визуализация корреляционных матриц. & Подключена как полезная библиотека для анализа, хотя в текущем варианте скрипта напрямую почти не используется. \\
\hline
\texttt{cluster} & Методы кластерного анализа. & Подключена для задач кластеризации; сам метод \texttt{kmeans()} доступен в базовом пакете \texttt{stats}. \\
\hline
\end{tabularx}
\end{table}

\subsection{Порядок подготовки файлов}

Перед запуском ЛР 2 необходимо проверить структуру проекта:

\begin{enumerate}[nosep]
  \item Открыть проект \texttt{r\_studio\_labs.Rproj} или папку проекта в RStudio.
  \item Убедиться, что файл \texttt{lab2\_script1.R} находится в папке \texttt{lab2\_r\_analyse}.
  \item Убедиться, что файл \texttt{lab2\_script2.R} находится в той же папке.
  \item Запустить сначала \texttt{lab2\_script1.R}. После его выполнения должен появиться файл \texttt{users\_data.csv}.
  \item Затем запустить \texttt{lab2\_script2.R}. Он читает \texttt{users\_data.csv}, добавляет кластеры и сохраняет графики \texttt{cluster\_scatter.png} и \texttt{cluster\_bar.png}.
\end{enumerate}

Ключевой момент: второй скрипт зависит от результата первого. Если файл \texttt{users\_data.csv} не создан или лежит не в рабочей папке, команда \texttt{read.csv()} выдаст ошибку.

\newpage

% =========================================================
% РАЗДЕЛ 5
% =========================================================
\section{ЛАБОРАТОРНАЯ РАБОТА №2: МЕТОДЫ, ПАРАМЕТРЫ И ВЛИЯНИЕ НА РАБОТУ}

\subsection{Скрипт 1. Подключение к API и сохранение данных}

В начале первого скрипта подключаются библиотеки:

\begin{lstlisting}
library(httr)
library(jsonlite)
\end{lstlisting}

Метод \texttt{library()} загружает установленный пакет в текущий сеанс R. Если пакет установлен, но не подключён, его функции могут быть недоступны по короткому имени.

\begin{table}[H]
\caption{Метод \texttt{library()}}
\begin{tabularx}{\textwidth}{|l|Y|}
\hline
\textbf{Параметр} & \textbf{Назначение и влияние} \\
\hline
Имя пакета & Указывает, какой пакет подключить. Например, \texttt{library(httr)} делает доступной функцию \texttt{GET()}. \\
\hline
\end{tabularx}
\end{table}

Далее указывается адрес удалённого ресурса и выполняется GET-запрос:

\begin{lstlisting}
url <- "https://jsonplaceholder.typicode.com/users"
response <- GET(url)
\end{lstlisting}

Метод \texttt{GET()} из библиотеки \texttt{httr} отправляет HTTP-запрос на сервер. В лабораторной он используется для получения списка пользователей из тестового API JSONPlaceholder.

\begin{table}[H]
\caption{Метод \texttt{GET()}}
\begin{tabularx}{\textwidth}{|l|Y|}
\hline
\textbf{Параметр} & \textbf{Назначение и влияние} \\
\hline
\texttt{url} & Адрес ресурса. Если адрес указан неверно, сервер вернёт ошибку или запрос не выполнится. \\
\hline
\texttt{timeout()} & Рекомендуемый дополнительный параметр для реальных проектов. Ограничивает время ожидания ответа, чтобы программа не зависала слишком долго. \\
\hline
\texttt{add\_headers()} & Используется, если API требует заголовки, например токен авторизации. В JSONPlaceholder авторизация не нужна. \\
\hline
\end{tabularx}
\end{table}

Ответ сервера проверяется по HTTP-коду:

\begin{lstlisting}
if (status_code(response) == 200) {
  cat("Подключение успешно. Код ответа:", status_code(response), "\n")
} else {
  stop("Ошибка подключения: ", status_code(response))
}
\end{lstlisting}

Метод \texttt{status\_code()} возвращает код ответа сервера. Код \texttt{200} означает успешный запрос. Метод \texttt{cat()} выводит понятное сообщение в консоль, а \texttt{stop()} останавливает выполнение программы при ошибке.

После проверки ответ преобразуется в текст, а затем в таблицу:

\begin{lstlisting}
users_json <- content(response, as = "text", encoding = "UTF-8")
users_data <- fromJSON(users_json, flatten = TRUE)
\end{lstlisting}

\begin{table}[H]
\caption{Получение и разбор JSON}
\begin{tabularx}{\textwidth}{|l|Y|Y|}
\hline
\textbf{Метод} & \textbf{Важные параметры} & \textbf{Как влияет на результат} \\
\hline
\texttt{content()} & \texttt{as = "text"} & Ответ сервера возвращается как обычная строка JSON, которую удобно передать в \texttt{fromJSON()}. \\
\hline
\texttt{content()} & \texttt{encoding = "UTF-8"} & Корректно обрабатывает символы в кодировке UTF-8. Это особенно важно для текстовых данных. \\
\hline
\texttt{fromJSON()} & \texttt{flatten = TRUE} & Вложенные поля JSON разворачиваются в плоские столбцы, например \texttt{address.geo.lat}. Без этого вложенные структуры было бы сложнее анализировать. \\
\hline
\end{tabularx}
\end{table}

В конце данные сохраняются в CSV-файл:

\begin{lstlisting}
write.csv(users_data, file = "users_data.csv",
          row.names = FALSE, fileEncoding = "UTF-8")
\end{lstlisting}

Метод \texttt{write.csv()} сохраняет таблицу R в файл CSV. Такой файл можно открыть не только в R, но и в табличном редакторе.

\begin{table}[H]
\caption{Параметры \texttt{write.csv()}}
\begin{tabularx}{\textwidth}{|l|Y|}
\hline
\textbf{Параметр} & \textbf{Назначение и влияние} \\
\hline
\texttt{users\_data} & Таблица, которую нужно сохранить. \\
\hline
\texttt{file = "users\_data.csv"} & Имя выходного файла. Файл создаётся в текущей рабочей папке. \\
\hline
\texttt{row.names = FALSE} & Не записывает отдельный столбец с номерами строк. Для учебного CSV это удобнее и чище. \\
\hline
\texttt{fileEncoding = "UTF-8"} & Сохраняет файл в кодировке UTF-8, чтобы текст корректно читался на разных системах. \\
\hline
\end{tabularx}
\end{table}

\subsection{Скрипт 2. Загрузка, анализ и визуализация}

Второй скрипт начинается с чтения CSV-файла:

\begin{lstlisting}
users_data <- read.csv("users_data.csv", fileEncoding = "UTF-8")
\end{lstlisting}

Метод \texttt{read.csv()} загружает CSV-файл в \texttt{data.frame}. Параметр \texttt{fileEncoding = "UTF-8"} помогает избежать проблем с кодировкой.

Затем выбираются координаты пользователей:

\begin{lstlisting}
numeric_data <- users_data %>%
  select(address.geo.lat, address.geo.lng) %>%
  mutate(
    address.geo.lat = as.numeric(address.geo.lat),
    address.geo.lng = as.numeric(address.geo.lng)
  )
\end{lstlisting}

Оператор \texttt{\%>\%} из экосистемы \texttt{dplyr} читается как <<передать результат дальше>>. Он позволяет записывать обработку данных цепочкой: сначала берём таблицу, затем выбираем нужные столбцы, затем преобразуем типы.

\begin{table}[H]
\caption{Подготовка числовых признаков}
\begin{tabularx}{\textwidth}{|l|Y|Y|}
\hline
\textbf{Метод} & \textbf{Назначение} & \textbf{Важные параметры} \\
\hline
\texttt{select()} & Оставляет только нужные столбцы. & Указываются имена столбцов: \texttt{address.geo.lat}, \texttt{address.geo.lng}. \\
\hline
\texttt{mutate()} & Создаёт новые столбцы или изменяет существующие. & В примере координаты перезаписываются после преобразования к числу. \\
\hline
\texttt{as.numeric()} & Преобразует значение в числовой тип. & Без числового типа \texttt{kmeans()} не сможет корректно посчитать расстояния между точками. \\
\hline
\end{tabularx}
\end{table}

Кластеризация выполняется методом k-средних:

\begin{lstlisting}
set.seed(42)
kmeans_result <- kmeans(numeric_data, centers = 3, nstart = 25)
users_data$cluster <- as.factor(kmeans_result$cluster)
\end{lstlisting}

Метод \texttt{kmeans()} группирует объекты так, чтобы точки внутри одного кластера были ближе друг к другу, чем к точкам других кластеров. В данной работе пользователи группируются по широте и долготе.

\begin{table}[H]
\caption{Параметры кластеризации}
\begin{tabularx}{\textwidth}{|l|Y|}
\hline
\textbf{Метод или параметр} & \textbf{Назначение и влияние} \\
\hline
\texttt{set.seed(42)} & Фиксирует генератор случайных чисел. Это нужно, чтобы при повторном запуске получались одинаковые кластеры. Число \texttt{42} можно заменить, но тогда результат может немного измениться. \\
\hline
\texttt{kmeans()} & Запускает алгоритм k-средних для числовых данных. \\
\hline
\texttt{centers = 3} & Количество кластеров. В лабораторной выбрано 3: этого достаточно, чтобы разделить 10 пользователей на несколько групп, но не дробить данные слишком мелко. \\
\hline
\texttt{nstart = 25} & Количество случайных начальных запусков алгоритма. Чем больше значение, тем выше шанс найти более удачное разбиение. Для небольшой выборки \texttt{25} является хорошим практическим значением. \\
\hline
\texttt{as.factor()} & Превращает номер кластера в категорию. Это удобно для раскраски графиков в \texttt{ggplot2}. \\
\hline
\end{tabularx}
\end{table}

Для визуализации используется \texttt{ggplot2}:

\begin{lstlisting}
p1 <- ggplot(users_data,
             aes(x = as.numeric(address.geo.lng),
                 y = as.numeric(address.geo.lat),
                 color = cluster)) +
  geom_point(size = 4) +
  geom_text(aes(label = name), vjust = -1, size = 3) +
  labs(title = "Кластеризация пользователей по геолокации",
       x = "Долгота (lng)",
       y = "Широта (lat)",
       color = "Кластер") +
  theme_minimal()
\end{lstlisting}

Метод \texttt{ggplot()} создаёт основу графика, а остальные методы добавляют слои и оформление.

\begin{table}[H]
\caption{Методы \texttt{ggplot2} для диаграммы рассеяния}
\begin{tabularx}{\textwidth}{|l|Y|Y|}
\hline
\textbf{Метод} & \textbf{Для чего используется} & \textbf{Важные параметры} \\
\hline
\texttt{ggplot()} & Создаёт объект графика. & Первый аргумент -- таблица данных. \\
\hline
\texttt{aes()} & Задаёт соответствия между данными и внешним видом графика. & \texttt{x} -- долгота, \texttt{y} -- широта, \texttt{color = cluster} раскрашивает точки по кластерам. \\
\hline
\texttt{geom\_point()} & Добавляет точки на график. & \texttt{size = 4} задаёт размер точек. \\
\hline
\texttt{geom\_text()} & Добавляет подписи к точкам. & \texttt{label = name} подписывает точки именами пользователей, \texttt{vjust = -1} поднимает подпись выше точки. \\
\hline
\texttt{labs()} & Задаёт заголовок, подписи осей и легенды. & Чем понятнее подписи, тем легче читать результат. \\
\hline
\texttt{theme\_minimal()} & Применяет аккуратную минималистичную тему. & Убирает лишнюю визуальную нагрузку. \\
\hline
\end{tabularx}
\end{table}

Для распределения пользователей по кластерам строится столбиковая диаграмма:

\begin{lstlisting}
p2 <- ggplot(users_data, aes(x = cluster, fill = cluster)) +
  geom_bar() +
  labs(title = "Распределение пользователей по кластерам",
       x = "Кластер",
       y = "Количество пользователей") +
  theme_minimal()
\end{lstlisting}

Метод \texttt{geom\_bar()} сам подсчитывает количество наблюдений в каждой категории \texttt{x}. Поэтому отдельно считать число пользователей по кластерам не требуется.

В конце графики сохраняются:

\begin{lstlisting}
ggsave("cluster_scatter.png", plot = p1, width = 10, height = 7, dpi = 150)
ggsave("cluster_bar.png", plot = p2, width = 8, height = 6, dpi = 150)
\end{lstlisting}

Метод \texttt{ggsave()} сохраняет последний или явно указанный график в файл.

\begin{table}[H]
\caption{Параметры \texttt{ggsave()}}
\begin{tabularx}{\textwidth}{|l|Y|}
\hline
\textbf{Параметр} & \textbf{Назначение и влияние} \\
\hline
Имя файла & Расширение определяет формат: \texttt{.png}, \texttt{.pdf}, \texttt{.jpg}. \\
\hline
\texttt{plot} & Указывает, какой график сохранить. Без этого параметра сохраняется последний выведенный график. \\
\hline
\texttt{width}, \texttt{height} & Размер изображения. Чем больше значения, тем крупнее итоговый файл и свободнее расположение подписей. \\
\hline
\texttt{dpi = 150} & Разрешение изображения. Большее значение даёт более чёткую картинку, но увеличивает размер файла. \\
\hline
\end{tabularx}
\end{table}

\newpage

% =========================================================
% РАЗДЕЛ 6
% =========================================================
\section{ПРАКТИЧЕСКИЕ ПОДСКАЗКИ И ТИПОВЫЕ ОШИБКИ}

\begin{table}[H]
\caption{Типовые проблемы при выполнении ЛР 1 и ЛР 2}
\begin{tabularx}{\textwidth}{|Y|Y|Y|}
\hline
\textbf{Проблема} & \textbf{Почему возникает} & \textbf{Как исправить} \\
\hline
\texttt{there is no package called ...} & Пакет не установлен. & Выполнить \texttt{install.packages("имя\_пакета")}, затем снова \texttt{library()}. \\
\hline
\texttt{could not find function ...} & Пакет установлен, но не подключён, либо в имени функции опечатка. & Проверить строку \texttt{library(...)} и имя метода. \\
\hline
\texttt{cannot open file 'users\_data.csv'} & Второй скрипт запущен раньше первого или рабочая папка выбрана неверно. & Сначала запустить \texttt{lab2\_script1.R}, затем проверить файл во вкладке \texttt{Files}. \\
\hline
График не появился & Выполнена не вся команда или график построен, но вкладка \texttt{Plots} скрыта. & Перейти во вкладку \texttt{Plots}, выполнить команду \texttt{print(p1)} или \texttt{print(p2)}. \\
\hline
\texttt{kmeans()} выдаёт ошибку по типам данных & Координаты прочитались как текст, а не как числа. & Использовать \texttt{as.numeric()} для столбцов широты и долготы. \\
\hline
CSV создался не там, где ожидалось & Текущая рабочая папка отличается от папки со скриптом. & Проверить \texttt{getwd()} или открыть проект через файл \texttt{.Rproj}. \\
\hline
\end{tabularx}
\end{table}

Перед сдачей удобно пройти короткий чек-лист:

\begin{enumerate}[nosep]
  \item В \texttt{Console} нет красных сообщений об ошибках.
  \item Во вкладке \texttt{Environment} видны основные объекты: \texttt{data1}, \texttt{df}, \texttt{users\_data}, \texttt{kmeans\_result}.
  \item Во вкладке \texttt{Files} есть \texttt{users\_data.csv}, \texttt{cluster\_scatter.png}, \texttt{cluster\_bar.png}.
  \item Во вкладке \texttt{Plots} отображаются построенные графики.
  \item В отчёте указано, какие методы применялись, какие параметры были выбраны и почему.
\end{enumerate}

\newpage

% =========================================================
% ЗАКЛЮЧЕНИЕ
% =========================================================
\addcontentsline{toc}{section}{Заключение}
\section*{ЗАКЛЮЧЕНИЕ}

В результате выполнения лабораторных работ №1 и №2 студент получает не только набор готовых R-скриптов, но и понимание процесса работы в RStudio. В ЛР 1 осваиваются базовые структуры данных языка R и встроенные средства визуализации. В ЛР 2 к ним добавляются внешние библиотеки, загрузка данных из API, сохранение CSV-файла, кластеризация методом k-средних и построение графиков через \texttt{ggplot2}.

Главная практическая идея такая: любой скрипт лучше выполнять не вслепую, а по шагам. После каждой важной команды надо смотреть, что изменилось в \texttt{Console}, \texttt{Environment}, \texttt{Files} и \texttt{Plots}. Тогда лабораторная превращается не в переписывание кода, а в понятный процесс анализа данных.

\newpage

% =========================================================
% СПИСОК ИСПОЛЬЗОВАННЫХ ИСТОЧНИКОВ
% =========================================================
\addcontentsline{toc}{section}{Список использованных источников}
\section*{СПИСОК ИСПОЛЬЗОВАННЫХ ИСТОЧНИКОВ}

\begin{enumerate}[label=\arabic*,leftmargin=1.25cm]
  \item The Comprehensive R Archive Network [Электронный ресурс].~-- Режим доступа: \url{https://cran.r-project.org/}.~-- Дата доступа: 15.06.2026.
  \item Posit. RStudio Desktop [Электронный ресурс].~-- Режим доступа: \url{https://posit.co/downloads/}.~-- Дата доступа: 15.06.2026.
  \item RStudio IDE User Guide [Электронный ресурс].~-- Режим доступа: \url{https://docs.posit.co/ide/user/}.~-- Дата доступа: 15.06.2026.
  \item JSONPlaceholder~-- Free Fake REST API [Электронный ресурс].~-- Режим доступа: \url{https://jsonplaceholder.typicode.com/}.~-- Дата доступа: 15.06.2026.
  \item Wickham,~H. ggplot2: Elegant Graphics for Data Analysis / H.~Wickham.~-- 2nd ed.~-- New York~: Springer, 2016.~-- 260~p.
  \item Wickham,~H. R for Data Science / H.~Wickham, G.~Grolemund.~-- Sebastopol~: O'Reilly Media, 2017.~-- 492~p.
\end{enumerate}

\end{document}
